\documentclass[12pt,spanish,mexico]{article}
\usepackage[utf8]{inputenc}
\PassOptionsToPackage{hyphens}{url}
\usepackage{hyperref}
\usepackage[es-nodecimaldot]{babel} % nombre de opción corregido
\usepackage{amsthm, amssymb, amsmath, amsfonts, bbm, mathtools}
\usepackage{csquotes}
\usepackage{enumerate}
\usepackage[shortlabels]{enumitem}
\usepackage{titling}
\usepackage{blindtext}
\usepackage{lmodern} 
\usepackage{float}
\usepackage{mathrsfs}
\usepackage[capitalise]{cleveref}
\usepackage{caption} 
\usepackage{graphicx}
\usepackage{chngcntr}
\usepackage{enumitem}    
\usepackage{cleveref} 
\usepackage[ruled,vlined]{algorithm2e}

\counterwithin{figure}{section}
\numberwithin{equation}{section}
\usepackage[vmargin=1in, hmargin=0.5in]{geometry}  

\setlength{\topmargin}{-50pt}
\setlength{\headheight}{38.23112pt}
\setlength{\headsep}{24pt}

\makeatletter
\def\iftitle{\ifx\thetitle \empty \else }
\def\safetitle{\iftitle{\thetitle}\fi}
\def\headertitle{%
    \ifdefined\hw{\hw}\fi%
    \ifdefined\hw \iftitle{ -- }\fi%
    \safetitle%
}
\def\@maketitle{%
  \begin{center}%
Proyecto: Asimilación de datos 4D-Var con el modelo logístico \\
Muñoz Macías Christian Jaime 
    \vskip 1em%
    \hrule
  \end{center}%
  \par
  \vskip 1.5em%
}
\makeatother
\usepackage{amsmath}
\usepackage{tikz}

\newcommand\RADot{%
  \mathord{%
    \mspace{1mu}%
    \text{\radot}%
    \mspace{1mu}%
  }%
}

\newcommand\radot{%
    \tikz[line cap=round,x=1ex,y=1ex,line width=0.3pt]
    {\draw (0,1) |- (1,0) (0.55,0) arc(0:90:0.55); \fill (0.23,0.23) circle (0.05);}%
}

%% encabezado personalizado
\usepackage{fancyhdr}
%% estilo plain personalizado
\fancypagestyle{plain}{%
    \fancyhf{} % limpia encabezado y pie de página
    \renewcommand{\headrulewidth}{0pt}
    \setlength{\headheight}{0pt}
    \fancyfoot[R]{\thepage}
}
%% estilo para todas las páginas
\pagestyle{fancy}
\fancyhf{}
\fancyhead[C]{Maestría en Matemáticas Aplicadas}
\fancyfoot[R]{\thepage}

%%%%%%%%%%%%%%%%%%%%%%%%
%% DEFINICIONES ÚTILES %%
%%%%%%%%%%%%%%%%%%%%%%%%
%% definiciones matemáticas
\newcommand{\N}{\mathbb{N}}
\newcommand{\Z}{\mathbb{Z}}
\newcommand{\Ha}{\mathcal{H}}
\newcommand{\R}{\mathbb{R}}
\newcommand{\Q}{\mathbb{Q}}
\newcommand{\C}{\mathbb{C}}
\newcommand{\PP}{\mathbb{P}}

\newtheorem{proposition}{Proposición}
\newtheorem{theorem}{Teorema}
\newtheorem{lemma}{Lema}
\newtheorem{corollary}{Corolario}

\theoremstyle{definition}
\newtheorem{problem}{Problema}
\newtheorem{definition}{Definición}[section]
\newtheorem*{example}{Ejemplo}

\theoremstyle{remark}
\newtheorem*{observation}{Observación}
\newtheorem*{note}{Nota}

\newenvironment{solution}{\begin{proof}[Solución]}{\end{proof}}
\newenvironment{dem}{\begin{proof}[Demostración]}{\end{proof}}
\renewcommand{\qedsymbol}{$\blacksquare$}

%% comandos útiles
\newcommand{\eps}{\varepsilon}
\newcommand{\To}{\longrightarrow}
\newcommand{\Eu}{\varPhi}
\newcommand{\Mapsto}{\longmapsto}

%% delimitadores emparejados
\DeclarePairedDelimiter{\set}{\{}{\}}
\DeclarePairedDelimiter{\floor}\lfloor\rfloor
\DeclarePairedDelimiter{\ceil}\lceil\rceil
\DeclarePairedDelimiter{\abs}\lvert\rvert
\DeclarePairedDelimiter{\gen}\langle\rangle
\DeclareMathOperator{\sgn}{sgn}
\DeclareMathOperator{\ord}{ord}

%% operador
\newcommand{\op}[1]{\operatorname{#1}}

%% etiquetas del documento
\title{Proyecto: Asimilación de datos 4D-Var con el modelo logístico}
\def\course{Materia}
\def\university{CIMAT - MsC. Maestría en matemáticas aplicadas}
\author{Muñoz Macías Christian Jaime}
\def\email{\url{christian.munoz@cimat.mx}}
\date{\today}

%% permitir saltos de línea en ecuaciones
\allowdisplaybreaks[4]

\begin{document}

\maketitle

\section{Modelo dinámico}
La dinámica del estado \(x(t)\) está gobernada por la ecuación logística
\[
\frac{dx}{dt} = r x \left(1 - \frac{x}{K}\right), \quad t \in [0, T], \qquad (1)
\]
con condición inicial
\[
x(0) = X_0, \qquad (2)
\]
donde los parámetros \(r > 0\) y \(K > 0\) se suponen conocidos y la incógnita es el estado inicial \(X_0\). Para hacer explícita la dependencia del control, denotamos por \(x(t; X_0)\) la solución de (1)–(2), y abreviamos \(x^\dagger(t) := x(t; X_0^\dagger)\) para la trayectoria verdadera. Disponemos de \(N\) observaciones ruidosas
\[
y_i = x^\dagger(t_i) + \eta_i, \quad \eta_i \sim \mathcal{N}(0, \sigma^2), \quad i = 1, \dots, N, \qquad (3)
\]
en tiempos \(0 < t_1 < \dots < t_N \leq T\) distribuidos a lo largo de la ventana de asimilación \([0, T]\). Adicionalmente contamos con un valor de background \(X_0^b\) —una estimación previa de la condición inicial— con incertidumbre caracterizada por la varianza \(B > 0\).

\section{Funcional 4D-Var}
El método 4D-Var formula la asimilación como el problema de optimización
\[
\min_{X_0 \in \mathbb{R}} J(X_0) = \underbrace{\frac{1}{2B}(X_0 - X_0^b)^2}_{\text{término de background } J_b} + \underbrace{\frac{1}{2\sigma^2} \sum_{i=1}^N (x(t_i; X_0) - y_i)^2}_{\text{término de observaciones } J_o} \qquad (4)
\]
sujeto a la restricción dinámica (1)–(2). La solución \(\hat{X}_0\) de (4) es el análisis; la trayectoria \(x(t; \hat{X}_0)\) obtenida al integrar el modelo desde \(\hat{X}_0\) es el análisis 4D. Bayesianamente, \(\hat{X}_0\) es el estimador de máximo a posteriori bajo verosimilitud gaussiana de los errores y prior gaussiana \(X_0 \sim \mathcal{N}(X_0^b, B)\).

\section{Datos del problema}
Los parámetros son como se indica en la siguiente tabla:

\begin{table}[H]
\centering
\begin{tabular}{|l|l|c|}
\hline
Nombre & Significado & Valor \\
\hline
\(r\) & Tasa de crecimiento (conocida) & 1.0 \\
\(K\) & Capacidad de carga (conocida) & 100 \\
\(T\) & Horizonte / ventana de asimilación & 5 \\
\(X_0^{\dagger}\) & Estado inicial verdadero (oculto) & 10 \\
\(X_0^{b}\) & Estado inicial background & por elegir \\
\(B\) & Varianza del background & por elegir \\
\(N\) & Número de observaciones en \([0, T]\) & 10 \\
\(\sigma\) & Desviación estándar del ruido & \(0.05 \cdot \max_{t\in[0,T]} x^{\dagger}(t)\) \\
\(\tau_k\) & Paso de descenso en iteración \(k\) & búsqueda de línea \\
\hline
\end{tabular}
\end{table}

\textbf{Observación importante sobre la ventana.} Se sugiere \(T = 5\). Con \(r = 1\) y \(X_0^{\dagger} = 10\), el sistema satura en \(t \approx 7\); tomar la ventana hasta \(t = 5\) asegura que las observaciones contengan información sobre la fase de crecimiento, donde \(X_0\) es realmente identificable. 

\textbf{Discusión: efecto del tamaño de la ventana de asimilación.}
En la observación anterior se plantean dos casos: cuando $x \approx K$ hay
saturación, y cuando el término $\frac{x(t)}{K} \approx 0$ domina el
comportamiento exponencial. Analizando estos dos casos:

\begin{itemize}
    \item \textbf{$T$ demasiado grande (saturación):} si la ventana se
    extiende hasta la zona plana de la logística, la mayoría de las
    observaciones caen en la región donde $x(t) \approx K$ y la
    trayectoria es prácticamente insensible a $X_0$. Con ruido gaussiano
    superpuesto, el funcional $J_o$ queda dominado por esas observaciones
    no informativas, y el optimizador ajusta el modelo a la fase saturada
    sin poder distinguir el valor correcto de $X_0$. El análisis resultante
    puede ser muy lejano al valor verdadero.

    \item \textbf{$T$ demasiado pequeño (poca información):} si la ventana
    abarca únicamente el inicio de la trayectoria, las observaciones cubren
    solo la fase de crecimiento casi exponencial, en la cual la logística
    es localmente indistinguible de $x(t) \approx X_0 e^{rt}$. En ese
    régimen, $K$ apenas influye en la solución y la información disponible
    es insuficiente para caracterizar el comportamiento completo del modelo,
    lo que produce una estimación de $X_0$ igualmente poco confiable.
\end{itemize}


\section{Resolución del problema de asimilación 4D-Var}

\subsection{Deducción de la ecuación adjunta}
Partiendo del Lagrangiano (5), calcula la variación \(\langle \partial\mathcal{L}/\partial x, h \rangle\) en una dirección arbitraria \(h(t)\) con \(h(0)\) libre, integra por partes y anula la variación. Obtén la ecuación adjunta
\[
-\dot{\lambda} - r \left( 1 - \frac{2x(t; X_0)}{K} \right) \lambda = -\frac{1}{\sigma^2} \sum_{i=1}^{N} \left( x(t_i; X_0) - y_i \right) \delta(t - t_i), \quad \lambda(T) = 0, \quad (6)
\]
junto con la condición de transversalidad \(\mu = \lambda(0)\) que relaciona el multiplicador de la EDO con el de la condición inicial. Aquí \(x(t; X_0)\) es la solución forward para el valor actual del control, es decir, el coeficiente de la EDO adjunta depende de la trayectoria forward y debe almacenarse durante la simulación del problema directo.

\begin{solution}
El funcional es:
\[
J(X_0) = \frac{1}{2B}(X_0 - X_0^b)^2 + \frac{1}{2\sigma^2} \sum_{i=1}^N (x(t_i; X_0) - y_i)^2
\]
Tenemos como Lagrangiano:
\[
\mathcal{L}(x, X_0, \lambda) = J(X_0) + \int_0^T \lambda(t) \left[ \dot{x} - r x(1 - x/K) \right] dt + \mu(x(0) - X_0), 
\]
sabemos que minimizar tiene que cumplir que: 
$\nabla J= \nabla \mathcal{L}$
entonces la derivada direccional de $J$ es:
\begin{align*}
D J(X_0)h =& D \mathcal{L}(x,X_0)(\xi,h) \\
=& D_x \mathcal{L} \xi+D_{X_0} \mathcal{L} h
\end{align*}


Lo anterior es válido porque en el punto estacionario las variaciones respecto a $\lambda$ se anulan (recuperan la ecuación forward), de modo que $D\mathcal{L}$ colapsa a $DJ$.

derivando respecto a $x$:
\[
D_x \mathcal{L} \xi = \frac{1}{\sigma^2} \sum_{i=1}^N (x(t_i) - y_i) \xi(t_i)+  \int_0^T \lambda(t) \left[ \dot{\xi}-r\left(1 -\frac{2x}{K} \right)\xi \right] dt + \mu\,\xi(0),
\]

intercambiar el diferencial con la integral es válido por la regularidad de $\lambda$ y $\xi$.

derivando respecto a $X_0$:
\[
D_{X_0} \mathcal{L} h= \frac{1}{B}(X_0 - X_0^b)h - \mu h, 
\]

nos interesa que $D_x \mathcal{L} \xi=0$, para ello trabajamos el término:
\[
\int_0^T \lambda(t) \left[ \dot{\xi}-r\left(1 -\frac{2x}{K} \right)\xi \right] dt 
\]
usando la regla del producto en $\lambda(t)\dot{\xi}$:
\[
\frac{d}{dt}(\lambda \xi)=\dot{\lambda} \xi+\lambda \dot{\xi}
\implies
\lambda \dot{\xi}=\frac{d}{dt}(\lambda \xi)-\dot{\lambda} \xi
\]
entonces la integral queda:
\[
\int_0^T \frac{d}{dt}(\lambda \xi)-\dot{\lambda} \xi-r\left(1 -\frac{2x}{K} \right)\lambda \xi \;  dt 
\]
evaluando el término de frontera:
\[
\int_0^T \frac{d}{dt}(\lambda \xi)\, dt=\lambda(T)\xi(T)-\lambda(0)\xi(0)=-\lambda(0)\xi(0)
\]
donde usamos $\lambda(T)=0$. Sustituyendo en $D_x\mathcal{L}\,\xi = 0$:
\[
\frac{1}{\sigma^2} \sum_{i=1}^N (x(t_i) - y_i) \xi(t_i)+\int_0^T\left[-\dot{\lambda} -r\left(1 -\frac{2x}{K} \right)\lambda\right] \xi \;  dt + (\mu - \lambda(0))\,\xi(0)= 0, 
\]
como $\xi(0)$ es libre, para que esto se anule para toda dirección $\xi$ el coeficiente de $\xi(0)$ debe ser cero:
\[
\mu = \lambda(0) \qquad \text{(condición de transversalidad)}
\]
cancelado ese término, usando que $\xi(t_i)=\int_0^T \xi(t)\,\delta(t-t_i)\,dt$, la condición debe valer para todo $\xi$, lo que implica la igualdad puntual:
\[
-\dot{\lambda}-r\left(1 -\frac{2x}{K} \right)\lambda=- \frac{1}{\sigma^2} \sum_{i=1}^N (x(t_i) - y_i) \delta (t-t_i)
\]
hemos deducido la ecuación adjunta con condición terminal $\lambda(T)=0$.
\end{solution}
\subsection{Deducción del gradiente reducido}
Varía \(\mathcal{L}\) respecto a \(X_0\), usa que \(x, \lambda\) satisfacen las ecuaciones forward y adjunta, y deduce
\[
\frac{dJ}{dX_0} = \frac{X_0 - X_0^b}{B} - \lambda(0). \quad (7)
\]
Comenta el significado: el gradiente es la suma de un tirón hacia la estimación inicial (primer término) y la sensibilidad retrógrada de las observaciones al estado inicial (segundo término, \(\lambda(0)\)). Esto justifica llamar al adjunto el “transportador hacia atrás de la información de las observaciones”.

\begin{solution}
Variamos $\mathcal{L}$ respecto a $X_0$ en una dirección $h$:
\[
D_{X_0} \mathcal{L}\, h = \frac{1}{B}(X_0 - X_0^b)h - \mu h
\]
como $x$ y $\lambda$ satisfacen las ecuaciones forward y adjunta, esta variación coincide con $DJ(X_0)h$, entonces:
\[
\frac{dJ}{dX_0} h= \frac{1}{B}(X_0 - X_0^b)h - \mu h
\]
sustituyendo la condición de transversalidad $\mu = \lambda(0)$ obtenida en el punto anterior:
\[
\boxed{\frac{dJ}{dX_0} = \frac{X_0 - X_0^b}{B} - \lambda(0)}
\]
El primer término jala la solución hacia el background; el segundo, $-\lambda(0)$, transporta hacia atrás la información de las observaciones a través de la ecuación adjunta.
\\
\\
El gradiente tiene una interpretación clara: si se busca minimizar $J$, el término 
$\frac{X_0 - X_0^b}{B}$ empuja a $X_0$ hacia el background $X_0^b$, pues se anula 
exactamente cuando $X_0 = X_0^b$. Por otro lado, $\lambda(0)$ es la solución en $t=0$ 
de la ecuación adjunta (6), que es una EDO lineal de primer orden no homogénea retrógrada 
en $[0,T]$: su término forzante son las discrepancias $(x(t_i;X_0)-y_i)$ en los tiempos 
de observación, y al integrar de $T$ hacia $0$ propaga esa información hasta el estado 
inicial. De ahí que $-\lambda(0)$ mida la sensibilidad acumulada de todas las observaciones 
respecto a $X_0$, jalando el análisis hacia donde los datos son más consistentes.

\end{solution}

\subsection{Generación de datos sintéticos}
Adopta el protocolo estándar de validación 4D-Var:
\begin{enumerate}[label=(\alph*)]
\item Integra (1) con \(X_0^\dagger = 10\) usando RK4 sobre \([0, 5]\);

\item Muestra \(N = 10\) observaciones en tiempos equiespaciados \(t_i = 0.5\, i\) para \(i = 1, \dots, 10\);
\item Añade ruido gaussiano con \(\sigma = 0.05 \cdot \max_{t\in[0,T]} x^\dagger(t)\).
\end{enumerate}
La trayectoria verdadera te servirá únicamente como referencia para medir el error del análisis; el algoritmo nunca la verá.

\subsection{Implementación forward}
Implementa un integrador RK4 para (1) que reciba \(X_0\) y devuelva la trayectoria \(x(t; X_0)\) sobre una malla densa que contenga los \(t_i\). Almacena la trayectoria; el adjunto la necesitará como coeficiente.

\subsection{Implementación del adjunto}
Implementa la integración retrógrada de (6) de \(T\) a \(0\), partiendo de \(\lambda(T) = 0\). Al cruzar cada \(t_i\) de derecha a izquierda, incorpora el salto
\[
\lambda(t_i^-) = \lambda(t_i^+) + \frac{1}{\sigma^2}(x(t_i; X_0) - y_i).
\]
Verifica con cuidado la consistencia de signos. El signo del salto se obtiene integrando la \(\delta\) de Dirac de la ecuación adjunta (6) sobre un entorno \([t_i-\epsilon, t_i+\epsilon]\) y haciendo \(\epsilon \to 0\); con la convención de signos de (6) y la integración de \(T\) hacia \(0\), el incremento \(\lambda(t_i^-) - \lambda(t_i^+)\) es positivo cuando \(x(t_i; X_0) > y_i\). Un error de signo aquí es la causa más frecuente de fallo del test de Taylor (punto siguiente). Devuelve \(\lambda(0)\) y, con él, ensambla el gradiente (7).


\begin{figure}[H]
    \centering
    \includegraphics[width=0.5\linewidth]{C:/Users/chris/OneDrive/escritorio/python/CIMAT/problemas inversos/gra_x_lamb_B_0.1.pdf}
    \caption{Trayectoria verdadera $x^\dagger(t)$, trayectoria obtenida a partir del background $X_0^b$ y solución adjunta $\lambda(t)$, con $B = 0.1$.}
\end{figure}
\subsection{Test de Taylor}
Antes de optimizar, valida tu gradiente. Para una dirección $\delta X_0 \in \mathbb{R}$
arbitraria y para \\
$\varepsilon \in \{10^{-1}, 10^{-2}, \dots, 10^{-8}\}$, calcula
\[
E_1(\varepsilon) = |J(X_0 + \varepsilon \delta X_0) - J(X_0)|, \quad
E_2(\varepsilon) = |J(X_0 + \varepsilon \delta X_0) - (J(X_0) +
\varepsilon \tfrac{dJ}{dX_0} \delta X_0)|.
\]
$E_1$ debe decaer con pendiente 1 y $E_2$ con pendiente 2 en gráfica log-log,
ambos sobre al menos cuatro órdenes de magnitud. Si la pendiente de $E_2$ no
es 2.

\begin{figure}[H]
    \centering
    \includegraphics[width=0.5\linewidth]{C:/Users/chris/OneDrive/escritorio/python/CIMAT/problemas inversos/gra_E1_E2_B_0.1.pdf}
    \caption{Comparación de $E_1$ y $E_2$ en escala log-log para el test de Taylor.}
\end{figure}

La figura muestra el comportamiento de ambos indicadores en escala log-log.
Para $\varepsilon \in [10^{-8}, 10^{-5}]$ ambas curvas presentan pendiente~1,
es decir, en este régimen $E_1$ y $E_2$ se comportan de manera idéntica: los
errores de redondeo en aritmética de punto flotante dominan y enmascaran la
diferencia entre ambos indicadores. En el intervalo intermedio
$\varepsilon \in [10^{-5}, 10^{-4}]$ se aprecia una transición con pendiente
$1 < m < 2$, donde $E_2$ comienza a separarse de $E_1$. Finalmente, para
$\varepsilon \in [10^{-4}, 10^{-1}]$ la pendiente de $E_2$ alcanza el
valor~2 esperado, confirmando que en este régimen el término cuadrático
domina y el gradiente adjunto es consistente con el funcional. Se concluye
que la derivada calculada es correcta.

\subsection{Minimización 4D-Var}
Implementa descenso más pronunciado con búsqueda de línea de Armijo:
\begin{enumerate}[label=(\alph*)]
    \item Resolver el problema forward con $X_0^{(k)}$ para obtener $x_k(t)$;
    \item Resolver el problema adjunto para obtener $\lambda_k(0)$;
    \item Ensamblar el gradiente~(7);
    \item Aplicar búsqueda de línea de Armijo y actualizar
          $X_0^{(k+1)} = X_0^{(k)} - \tau_k \dfrac{dJ}{dX_0}$.
\end{enumerate}
Se inicializa con $X_0^{(0)} = X_0^b$ (el background). Se reporta la convergencia
de $J(X_0^{(k)})$, $|dJ/dX_0|$ y $|X_0^{(k)} - X_0^\dagger|$.

\subsection{Estudio del peso relativo background/observaciones}
Se realizan tres experimentos de asimilación variando la confianza relativa entre
el background y las observaciones, manteniendo $X_0^b = 5$ (alejado del valor
verdadero $X_0^\dagger = 10$):

\subsubsection*{(a) Background dominante: $B = 0.1$}

\begin{figure}[H]
    \centering
    \includegraphics[width=0.5\linewidth]{C:/Users/chris/OneDrive/escritorio/python/CIMAT/problemas inversos/gra_J_0.1.pdf}
    \caption{Evolución del funcional $J$ a lo largo de las iteraciones, $B = 0.1$.}
\end{figure}
\begin{figure}[H]
    \centering
    \includegraphics[width=0.5\linewidth]{C:/Users/chris/OneDrive/escritorio/python/CIMAT/problemas inversos/gra_nablaJ_B_0.1.pdf}
    \caption{Evolución de $|\nabla J|$ a lo largo de las iteraciones, $B = 0.1$.}
\end{figure}
\begin{figure}[H]
    \centering
    \includegraphics[width=0.5\linewidth]{C:/Users/chris/OneDrive/escritorio/python/CIMAT/problemas inversos/gra_X0_B_0.1.pdf}
    \caption{Evolución del error $|X_0^{(k)} - X_0^\dagger|$ a lo largo de
    las iteraciones, $B = 0.1$.}
\end{figure}

Con $B = 0.1$ el término de la estimación inicial domina ampliamente al término de
observaciones en el funcional $J$. El análisis converge a
$\hat{X}_0 \approx 6.2$, es decir, a menos de $1.2$ unidades del estimación inicial
$X_0^b = 5$ y a más de $3.8$ unidades del valor verdadero $X_0^\dagger = 10$.
Esta varianza tan pequeña penaliza fuertemente cualquier desviación respecto
al prior, de modo que las observaciones tienen poca capacidad para corregir
la estimación inicial. El resultado muestra un resultado conservador a la estimación inicial.

\subsubsection*{(b) Régimen equilibrado: $B = 10$}

\begin{figure}[H]
    \centering
    \includegraphics[width=0.5\linewidth]{C:/Users/chris/OneDrive/escritorio/python/CIMAT/problemas inversos/gra_J_10.0.pdf}
    \caption{Evolución del funcional $J$ a lo largo de las iteraciones, $B = 10$.}
\end{figure}
\begin{figure}[H]
    \centering
    \includegraphics[width=0.5\linewidth]{C:/Users/chris/OneDrive/escritorio/python/CIMAT/problemas inversos/gra_nablaJ_B_10.0.pdf}
    \caption{Evolución de $|\nabla J|$ a lo largo de las iteraciones, $B = 10$.}
\end{figure}
\begin{figure}[H]
    \centering
    \includegraphics[width=0.5\linewidth]{C:/Users/chris/OneDrive/escritorio/python/CIMAT/problemas inversos/gra_X0_B_10.0.pdf}
    \caption{Evolución del error $|X_0^{(k)} - X_0^\dagger|$ a lo largo de
    las iteraciones, $B = 10$.}
\end{figure}

Con $B = 10$ el peso relativo entre la estimación inicial y las observaciones es más
equilibrado. El análisis converge a $\hat{X}_0 \approx 9.35$, a solo $0.65$
unidades del valor verdadero $X_0^\dagger = 10$. Al otorgar mayor importancia
a los datos, el algoritmo puede alejarse significativamente del 
estado inicial y aproximarse al estado verdadero. Este caso representa un compromiso
razonable entre regularización y ajuste a las observaciones.


\subsubsection*{(c) Observaciones dominantes: $B = 10^4$}

\begin{figure}[H]
    \centering
    \includegraphics[width=0.5\linewidth]{C:/Users/chris/OneDrive/escritorio/python/CIMAT/problemas inversos/gra_J_10000.0.pdf}
    \caption{Evolución del funcional $J$ a lo largo de las iteraciones,
    $B = 10^4$.}
\end{figure}
\begin{figure}[H]
    \centering
    \includegraphics[width=0.5\linewidth]{C:/Users/chris/OneDrive/escritorio/python/CIMAT/problemas inversos/gra_nablaJ_B_10000.0.pdf}
    \caption{Evolución de $|\nabla J|$ a lo largo de las iteraciones,
    $B = 10^4$.}
\end{figure}
\begin{figure}[H]
    \centering
    \includegraphics[width=0.5\linewidth]{C:/Users/chris/OneDrive/escritorio/python/CIMAT/problemas inversos/gra_X0_B_10000.0.pdf}
    \caption{Evolución del error $|X_0^{(k)} - X_0^\dagger|$ a lo largo de
    las iteraciones, $B = 10^4$.}
\end{figure}

Con $B = 10^4$ el término de la estimación inicial se vuelve prácticamente inactivo y
el funcional está dominado por el ajuste a las observaciones. El análisis
converge a $\hat{X}_0 \approx 10.2$, el resultado más cercano al valor
verdadero $X_0^\dagger = 10$ de los tres experimentos, con un error absoluto
de apenas $0.2$ unidades. La ligera sobreestimación es atribuible al ruido
gaussiano en las observaciones, cuya realización particular eleva
marginalmente la media muestral. Este caso muestra que, cuando la estimación inicial
es poco confiable, confiar casi exclusivamente en los datos produce la mejor
recuperación del estado verdadero.

\bigskip
\noindent\textbf{Conclusión del estudio.}\quad
Los tres experimentos ilustran el papel regulador de $B$ en el funcional
4D-Var. Un valor pequeño de $B$ ancla la solución cerca de la estimación inicial e
impide que las observaciones corrijan estimaciones iniciales erróneas; un
valor grande relega el prior y permite que los datos guíen la optimización.
En la práctica, $B$ debe calibrarse de acuerdo con la incertidumbre real del
estimación inicial: sobreestimar esa confianza produce análisis sesgados hacia el
prior, mientras que subestimarla puede amplificar el ruido de las
observaciones. El caso $B = 10^4$ sugiere que, para este problema, la
información contenida en las diez observaciones es suficiente para recuperar
$X_0^\dagger$ con alta precisión incluso partiendo de una estimación inicial alejada de
$5$ unidades del valor verdadero.
\end{document}